\documentclass[a4paper,11pt]{article}
\usepackage[utf8]{inputenc}
\usepackage[T1]{fontenc}
\usepackage[french]{babel}
\usepackage[left=2cm,right=2cm,top=2cm,bottom=2cm]{geometry}
\usepackage{xcolor}
\usepackage{hyperref}
\usepackage{fontawesome5}
\usepackage{parskip}

% --- Couleurs Capgemini ---
\definecolor{primary}{RGB}{0, 70, 173}
\hypersetup{colorlinks=true, urlcolor=primary}

\begin{document}

% --- EXPÉDITEUR ---
\noindent
\begin{minipage}[t]{0.5\textwidth}
    \textbf{Loïc Aron MBASSI EWOLO}\\
    Élève-Ingénieur Bac+5 -- Double diplôme ENSEA\\[3pt]
    \faMapMarker\ Cergy, France\\
    \faPhone\ (+33) 7 59 52 33 98\\
    \faEnvelope\ \href{mailto:loicaron.mbassiewolo@ensea.fr}{loicaron.mbassiewolo@ensea.fr}
\end{minipage}
\hfill
% --- DATE ---
\begin{minipage}[t]{0.4\textwidth}
    \raggedleft
    Cergy, le 9 mars 2026
\end{minipage}

\vspace{1.2cm}

% --- DESTINATAIRE ---
\hfill
\begin{minipage}[t]{0.5\textwidth}
    \raggedleft
    \textbf{Capgemini}\\
    Service Recrutement\\
    Issy-les-Moulineaux, France
\end{minipage}

\vspace{1cm}

\noindent\textbf{Objet :} Candidature en alternance -- Ingénieur Logiciels Embarqués

\vspace{0.8cm}

Madame, Monsieur,

Développer en C/C++ sur Nvidia Jetson pour faire naviguer un robot autonome (SLAM, Lidar 3D, YOLOv10), puis concevoir des gants IoT sur STM32 fusionnant capteurs inertiels et vision pour traduire la langue des signes : ces deux projets m'ont ancré dans les contraintes du développement embarqué, du bas niveau hardware à l'optimisation logicielle sous contrainte temps réel. C'est cette expertise que je souhaite approfondir au sein du pôle Industrie et Développement Embarqué de Capgemini.

Le challenge CoHoMa (robotique autonome militaire, ENSEA/Armée française) m'a confronté au développement firmware multi-plateforme : déploiement C/C++ sur Jetson, intégration Lidar VLP16 et caméra de profondeur, containerisation Docker pour garantir la reproductibilité entre environnements. En parallèle, mon projet académique de simulation de flotte de taxis en C sous Linux (mémoire partagée, signaux POSIX, ordonnanceur FIFO) m'a donné une compréhension fine des mécanismes système : IPC, gestion manuelle de la mémoire, synchronisation de processus concurrents.

Mon stage au laboratoire IoT de l'École Polytechnique de Yaoundé a complété ce profil par l'optimisation de modèles ML pour microcontrôleurs (TinyML, TF Lite, quantification), sous contraintes mémoire strictes. À INRIA Saclay (équipe TRIBE, avril à août 2026), je développe actuellement des pipelines Python structurés et rédige en anglais dans un environnement de recherche exigeant, ce qui renforce mon autonomie technique et ma capacité à communiquer dans un contexte international.

Je suis disponible dès septembre 2026 pour une alternance de 12 mois. Contribuer à des projets firmware multi-plateforme chez Capgemini, aux côtés d'équipes intervenant sur des sujets comme les véhicules autonomes ou les énergies du futur, est une opportunité que je souhaite saisir.

Je reste à votre disposition pour un entretien afin de discuter de ma candidature.

Je vous prie d'agréer, Madame, Monsieur, l'expression de mes salutations distinguées.

\vspace{0.8cm}

\textbf{Loïc Aron MBASSI EWOLO}\\
\textit{Élève-Ingénieur ENSEA / Polytechnique Yaoundé}

\end{document}
